\chapter{机器学习项目流程}

\section{学习目标}

学完本章后，你应该能够：

\begin{itemize}
    \item 说清楚一个最小机器学习项目为什么必须从任务定义和数据切分开始，而不是从选模型开始；
    \item 区分训练集、验证集和测试集在项目中的不同角色；
    \item 理解为什么类别覆盖、分层切分和数据泄漏会直接决定评估是否可信；
    \item 用一个极小的恒星分类数据集走通“切分 $\rightarrow$ baseline $\rightarrow$ 训练 $\rightarrow$ 评估 $\rightarrow$ 误差检查”的完整骨架；
    \item 理解准确率、混淆矩阵和错误案例分析为什么比单个分数更重要。
\end{itemize}

\section{为什么流程先于模型}

在机器学习初学阶段，最常见的误区不是“模型太弱”，而是“流程根本没有建立起来”。很多人会本能地把注意力放在：

\begin{itemize}
    \item 该用 KNN 还是随机森林；
    \item 该不该上深度学习；
    \item 哪个库函数更高级。
\end{itemize}

但真正决定一个结果是否可信的，往往是更早的那些问题：

\begin{itemize}
    \item 任务到底定义清楚了吗；
    \item 标签是否可靠；
    \item 数据是否合理切分；
    \item baseline 是否存在；
    \item 测试分数是否独立于训练过程。
\end{itemize}

如果这些基础没有建立起来，模型即使复杂，也可能只是在放大前面流程中的错误。

\section{一个最小机器学习工作流}

对本科阶段最小可用的机器学习任务，我们可以把流程写成：

\begin{enumerate}
    \item 明确科学问题和预测目标；
    \item 定义样本单位、特征和标签；
    \item 先划分训练集、验证集和测试集；
    \item 建立一个清楚的 baseline；
    \item 训练一个数据驱动模型；
    \item 用独立数据评估结果；
    \item 回到错误案例和科学解释。
\end{enumerate}

本章最重要的思想是：\textbf{这些步骤并不是官僚流程，而是保证评估结果有意义的最小条件。}

\section{任务定义：先回答“要预测什么”}

任何工作流的第一步，都不是写代码，而是定义任务。例如在本章配套数据里，我们面对的是一个极小恒星分类问题：

\begin{itemize}
    \item 样本单位：一颗恒星；
    \item 特征：颜色指数和绝对星等；
    \item 标签：\texttt{main\_sequence}、\texttt{giant}、\texttt{white\_dwarf}；
    \item 目标：根据特征预测恒星类别。
\end{itemize}

只有把这几个对象说清楚，后面的切分、训练和评估才有共同语义。否则，你甚至可能在不知不觉中把“派生自标签的信息”当成输入特征，从而得到虚假的高分。

\section{为什么要先切分，再训练}

一个最基本但最容易被忽略的原则是：\textbf{数据切分必须发生在模型真正学习之前。}

理由很直接。如果一个模型在训练时已经间接看过测试样本，那么测试分数就不再代表泛化能力，只是在复述训练记忆。机器学习里把这种问题称为\textbf{数据泄漏（data leakage）}。

设总数据集为 $\mathcal{D}$，一个理想化的切分可以写成：

\[
\mathcal{D} = \mathcal{D}_{\mathrm{train}} \cup
\mathcal{D}_{\mathrm{val}} \cup
\mathcal{D}_{\mathrm{test}},
\]

并且三部分在评估用途上应尽量保持独立：

\begin{itemize}
    \item $\mathcal{D}_{\mathrm{train}}$：用于学习模型参数；
    \item $\mathcal{D}_{\mathrm{val}}$：用于比较模型和调参；
    \item $\mathcal{D}_{\mathrm{test}}$：只在最后用于报告独立评估结果。
\end{itemize}

本章的教学数据极小，因此 notebook 只演示训练集和测试集；验证集及更系统的模型选择，会放到后面章节展开。

\section{为什么“切分方式”本身就是评估设计}

很多初学者把切分看成一个技术细节，觉得“反正切一下就行”。但实际上，切分方式本身就是评估设计的一部分。

设想一种不合理的情况：把前几行全部当成训练集，后几行全部当成测试集。若数据恰好按类别排布，就会出现某类对象只在测试集中出现、训练集中完全没有的情况。此时模型并不是“学得不好”，而是根本没有机会学到那一类。

因此，好的切分至少应当避免两件事：

\begin{itemize}
    \item 测试集中出现训练集完全未见过的常规类别；
    \item 训练集与测试集在样本来源上存在明显不公平差异却未被说明。
\end{itemize}

\section{分层切分：为什么类别覆盖很重要}

在分类任务中，一个自然的要求是：训练集中最好每一类都有代表样本。对于小型教学数据，这一点尤其关键。

因此，本章 notebook 采用一个极简分层思路：每个类别各留出一个样本放到测试集，其余样本用于训练。代码很短，但教学含义非常强：

\begin{examplebox}[教学版分层切分]
\begin{lstlisting}[style=pythonstyle]
test_indices = {0, 4, 6}
train_rows = [row for index, row in enumerate(rows)
              if index not in test_indices]
test_rows = [row for index, row in enumerate(rows)
             if index in test_indices]
\end{lstlisting}
\end{examplebox}

这不是完整严谨的统计切分方案，但它足以让学生看到：\textbf{一个看似普通的索引选择，已经在深刻影响后面所有评估结果。}

\section{baseline：流程里的第一把尺子}

在真正训练模型前，我们总要先问：如果不用复杂学习器，只靠简单规则，能做到什么程度？

这正是 baseline 的作用。它不是一个“形式上必须有”的环节，而是判断复杂模型是否真的带来增益的第一把尺子。

在本章恒星分类例子里，baseline 可以是：

\begin{itemize}
    \item 来自前一章的规则分类器；
    \item 一个按类别中心做最近分类的最小数据驱动方法；
    \item 或者更简单的多数类猜测器。
\end{itemize}

只有先知道这些简单方法能做什么，我们才有资格判断一个更复杂模型是不是值得。

\section{训练一个最小分类器：最近类别中心}

本章 notebook 采用的最小数据驱动模型是\textbf{最近类别中心}分类器。它的思想是：

\begin{enumerate}
    \item 用训练集计算每一类在特征空间中的中心；
    \item 对一个新样本，计算它到各类中心的距离；
    \item 把它分到最近的那一类。
\end{enumerate}

这一步的教学意义，不在于模型强大，而在于它把“训练”具体化了：训练并不总是复杂优化，也可以只是根据训练集组织出类别结构。

\begin{examplebox}[构造最小分类器]
\begin{lstlisting}[style=pythonstyle]
def feature_vector(row):
    return (row["bp_rp"], row["absolute_g_mag"])

def build_centroids(training_rows):
    grouped = {}
    for row in training_rows:
        grouped.setdefault(row["stellar_class"], []).append(
            feature_vector(row)
        )
\end{lstlisting}
\end{examplebox}

\section{评估：为什么不能只说“看起来还行”}

一旦有了预测结果，我们需要用明确指标评估。对分类任务，最基本的度量之一是准确率：

\[
\mathrm{Accuracy} =
\frac{1}{N_{\mathrm{test}}}
\sum_{i=1}^{N_{\mathrm{test}}}
\mathbb{1}(\hat{y}_i = y_i),
\]

其中 $\mathbb{1}$ 表示指示函数，预测正确时取 $1$，否则取 $0$。

准确率直观，但信息并不完整。因为两个模型可能准确率相同，却犯了完全不同类型的错误：一个可能把所有白矮星都分错，另一个则只是偶尔混淆主序星和巨星。

\section{混淆矩阵：看清错误结构}

为了知道模型到底错在什么地方，我们常引入混淆矩阵。若真实类别为 $a$、预测类别为 $b$ 的样本数记为 $C_{ab}$，则混淆矩阵记录了所有“真实 $\rightarrow$ 预测”的对应次数。

这类结构能帮助我们回答：

\begin{itemize}
    \item 模型是不是系统性混淆某两类对象；
    \item 哪个类别最难识别；
    \item 误差是否来自类别不平衡或特征重叠。
\end{itemize}

\begin{examplebox}[一个最小混淆矩阵结构]
\begin{lstlisting}[style=pythonstyle]
{
    "main_sequence": {"main_sequence": 1, "giant": 0, "white_dwarf": 0},
    "giant": {"main_sequence": 1, "giant": 0, "white_dwarf": 0},
    "white_dwarf": {"main_sequence": 0, "giant": 0, "white_dwarf": 1},
}
\end{lstlisting}
\end{examplebox}

即使样本量很小，这种表示也比一句“准确率是 0.667”更有教学价值。

\section{错误案例分析：从分数回到样本}

真正重要的评估，不会停在分数上，而会回到具体样本：

\begin{itemize}
    \item 是哪个样本被错分了；
    \item 它的颜色和绝对星等落在什么区域；
    \item 错分是因为训练样本太少，还是因为两类本来就接近；
    \item 这个错误在科学上是否严重。
\end{itemize}

这一步非常像实验数据分析里的残差诊断。分数给你的是总体印象，错误案例才告诉你模型究竟学到了什么、没学到什么。

\section{小样本场景下为什么更要谨慎}

本章教学数据规模很小，因此所有结果都必须谨慎解释。小样本机器学习的典型风险包括：

\begin{itemize}
    \item 一次切分的偶然性会强烈影响分数；
    \item 某个类别只要少一个样本，边界就会显著变化；
    \item 几个样本的标签偏差就足以扭曲结论；
    \item 一个看起来很高的分数，可能只是因为任务被切分得太容易。
\end{itemize}

因此，这一章的目标不是“训练一个可靠模型”，而是建立一条清楚的工作流骨架。只有流程先可靠，后面更复杂的模型比较才有意义。

\section{为什么 notebook 还保留了 scikit-learn 对照}

配套 notebook 中有一个可选单元：如果环境安装了 \texttt{scikit-learn}，就可以再用标准库里的 KNN 分类器做对比。

这一步的作用不是“把你手写的代码替换掉”，而是帮助你理解两件事：

\begin{itemize}
    \item 自己写的最小流程，与标准机器学习库的核心思想是一致的；
    \item 成熟生态提供的是更规范的实现、更方便的接口和更广泛的模型选择，而不是新的魔法。
\end{itemize}

\section{可复现性：流程也需要被别人重走}

机器学习项目除了“能跑通”，还要“别人能重走”。因此，哪怕在小型教学例子里，也应该逐步养成以下习惯：

\begin{itemize}
    \item 明确写出切分方式；
    \item 记录使用了哪些特征；
    \item 明确 baseline 与正式模型的区别；
    \item 保留预测结果和错误样本列表；
    \item 说明环境依赖是必须项还是可选项。
\end{itemize}

这些工作看上去不像“模型创新”，却是把一个课堂示例变成科研工作流的关键。

\begin{warnbox}[机器学习流程中的典型问题]
\begin{itemize}
    \item 先看完整数据，再决定如何切分；
    \item 训练集没有覆盖全部类别，却仍然把测试分数当真；
    \item 只报告单一准确率，不展示错误样本；
    \item 在没有 baseline 的情况下直接上复杂模型；
    \item 用测试集反复改规则，直到分数看起来满意。
\end{itemize}
\end{warnbox}

\section{AI 助手如何帮助你}

在这一章，AI 助手适合帮助你：

\begin{itemize}
    \item 帮你检查一个切分方案是否明显不合理；
    \item 把最小分类流程整理成更清楚的函数结构；
    \item 帮你把准确率、混淆矩阵和错误案例分析写成更规范的评估报告；
    \item 提醒你哪些地方可能存在数据泄漏或类别覆盖不足。
\end{itemize}

但一个流程是否真的可信、测试集是否真的独立、以及样本规模是否足以支撑当前结论，仍然必须由你自己判断。

\section{练习}

\begin{exercisebox}[基础练习]
说明为什么“训练集中完全没有某一类样本”会让这个类别上的测试结果失去解释意义。
\end{exercisebox}

\begin{exercisebox}[进阶练习]
从准确率公式出发，解释为什么两个模型即使准确率相同，也仍然可能有完全不同的科学可用性。
\end{exercisebox}

\begin{exercisebox}[开放练习]
为你自己的一个机器学习小项目写一份流程草案，至少包含：任务定义、样本单位、特征、标签、切分策略、baseline、正式模型和评估指标。
\end{exercisebox}

\section{本章小结}

机器学习项目真正的起点，不是某个模型名称，而是一条可解释、可复现、可评估的流程。本章通过一个极小恒星分类例子，建立了从任务定义、数据切分和 baseline，到训练、评估和错误分析的最小骨架。只要这条骨架真正建立起来，后面进入回归、模型选择与更复杂算法时，你就不会把“会调库”误当成“会做机器学习项目”。
